\documentclass{article}
\usepackage[utf8]{inputenc}
\usepackage{graphicx}
\usepackage{grffile}
\usepackage{amsmath}
\usepackage{bm}
\usepackage{booktabs}
\usepackage{hyperref}

\title{Extension of the Hammer--Rosen Model to a Gray Two-Temperature System}
\author{Yuval Kochavi}
\date{\today}

\begin{document}

\maketitle

\section{Extension of the Hammer--Rosen Model to a Gray Two-Temperature System}

Hammer and Rosen considered the supersonic Marshak-wave problem under the
assumption of complete thermodynamic equilibrium, namely
\[
T_R = T_m \equiv T .
\]
In that case the radiation energy density is written as
\[
E = aT^4 ,
\]
and the material energy equation reduces to a single nonlinear diffusion
equation for the material temperature. In their notation, the material internal
energy and opacity are assumed to obey power laws of the form
\[
e = fT^\beta \rho^{-\mu},
\qquad
\frac{1}{\kappa} = gT^\alpha \rho^{-\lambda}.
\]
This leads to a nonlinear diffusion equation of the form
\[
\frac{\partial T^\beta}{\partial t}
=
C \frac{\partial^2 T^{4+\alpha}}{\partial x^2},
\]
where
\[
C =
\frac{16\sigma g}{3f(4+\alpha)}
\rho^{\mu-\lambda-2}.
\]
Hammer and Rosen then solve this equation perturbatively using the small
parameter
\[
\epsilon = \frac{\beta}{4+\alpha}.
\]
Their main result for the heat-front position is
\[
x_F^2(t)
=
\frac{2+\epsilon}{1-\epsilon}
C H^{-\epsilon}(t)
\int_0^t H(t')\,dt',
\]
where
\[
H(t) = T_s^{4+\alpha}(t).
\]
The absorbed energy per unit area is, to first order in \(\epsilon\),
\[
\mathcal{E}(t)
=
f\rho^{1-\mu}
\frac{x_F(t)H^\epsilon(t)}{1-\epsilon}.
\]
These are the one-temperature CTE results of Hammer and Rosen. :contentReference[oaicite:0]{index=0}

\section{Gray Two-Temperature Equations}

A more general description is obtained by allowing the radiation and material
temperatures to differ:
\[
T_R \neq T_m .
\]
The gray radiation-diffusion equations are
\begin{align}
\frac{\partial E}{\partial t}
-
\nabla \cdot
\left(
\frac{c}{3\sigma}
\nabla E
\right)
&=
c\sigma \left(aT_m^4 - E\right),
\\
\rho\frac{\partial e}{\partial t}
&=
c\sigma \left(E-aT_m^4\right),
\end{align}
where
\[
E = aT_R^4,
\qquad
\sigma = \kappa \rho .
\]

Using the same power-law assumptions as Hammer and Rosen,
\[
e = fT_m^\beta \rho^{-\mu},
\qquad
\frac{1}{\kappa} = gT_m^\alpha \rho^{-\lambda},
\]
we have
\[
\kappa = \frac{1}{g}T_m^{-\alpha}\rho^\lambda,
\]
and therefore
\[
\sigma = \kappa\rho
=
\frac{1}{g}T_m^{-\alpha}\rho^{1+\lambda}.
\]
Thus
\[
\frac{c}{3\sigma}
=
\frac{cg}{3}\rho^{-1-\lambda}T_m^\alpha .
\]

Defining
\[
D_0 = \frac{cg}{3}\rho^{-1-\lambda},
\qquad
S_0 = \frac{c}{g}\rho^{1+\lambda},
\qquad
A_m = f\rho^{1-\mu},
\]
the two-temperature gray equations become
\begin{align}
\frac{\partial E}{\partial t}
-
\nabla \cdot
\left(
D_0T_m^\alpha \nabla E
\right)
&=
S_0T_m^{-\alpha}
\left(aT_m^4-E\right),
\\
A_m\frac{\partial T_m^\beta}{\partial t}
&=
S_0T_m^{-\alpha}
\left(E-aT_m^4\right).
\end{align}

\section{Energy Conservation}

Adding the radiation and material equations cancels the emission-absorption
coupling term. Therefore,
\[
\frac{\partial}{\partial t}
\left(
E + A_mT_m^\beta
\right)
=
\nabla \cdot
\left(
D_0T_m^\alpha \nabla E
\right).
\]

In one spatial dimension this becomes
\[
\frac{\partial}{\partial t}
\left(
E + A_mT_m^\beta
\right)
=
\frac{\partial}{\partial x}
\left(
D_0T_m^\alpha
\frac{\partial E}{\partial x}
\right).
\]

The total areal energy inside the heated region is therefore
\[
\mathcal{E}(t)
=
\int_0^{x_F(t)}
\left[
E(x,t) + A_mT_m^\beta(x,t)
\right]dx .
\]

The incoming flux at the surface is
\[
F_{\rm in}(t)
=
-
\left.
D_0T_m^\alpha
\frac{\partial E}{\partial x}
\right|_{x=0}.
\]

Hence the exact energy-conservation relation is
\[
\boxed{
\frac{d\mathcal{E}}{dt}
=
F_{\rm in}(t)
}
\]
or explicitly,
\[
\boxed{
\frac{d}{dt}
\int_0^{x_F(t)}
\left[
E(x,t) + A_mT_m^\beta(x,t)
\right]dx
=
-
\left.
D_0T_m^\alpha
\frac{\partial E}{\partial x}
\right|_{x=0}.
}
\]

\section{Recovery of the Hammer--Rosen Limit}

The Hammer--Rosen equation is recovered when the coupling between radiation
and matter is sufficiently strong, so that
\[
E \simeq aT_m^4 .
\]
Equivalently,
\[
T_R \simeq T_m \equiv T .
\]

In this case, the total energy equation becomes
\[
\frac{\partial}{\partial t}
\left(
aT^4 + A_mT^\beta
\right)
=
\frac{\partial}{\partial x}
\left(
D_0T^\alpha
\frac{\partial aT^4}{\partial x}
\right).
\]

If the material energy dominates the radiation energy,
\[
A_mT^\beta \gg aT^4,
\]
then
\[
A_m\frac{\partial T^\beta}{\partial t}
=
\frac{\partial}{\partial x}
\left(
D_0T^\alpha
\frac{\partial aT^4}{\partial x}
\right).
\]

Since
\[
T^\alpha\frac{\partial T^4}{\partial x}
=
\frac{4}{4+\alpha}
\frac{\partial T^{4+\alpha}}{\partial x},
\]
we obtain
\[
A_m\frac{\partial T^\beta}{\partial t}
=
\frac{4aD_0}{4+\alpha}
\frac{\partial^2 T^{4+\alpha}}{\partial x^2}.
\]

Therefore,
\[
\frac{\partial T^\beta}{\partial t}
=
C
\frac{\partial^2 T^{4+\alpha}}{\partial x^2},
\]
with
\[
\boxed{
C =
\frac{4aD_0}{(4+\alpha)A_m}
=
\frac{4acg}{3(4+\alpha)f}
\rho^{\mu-\lambda-2}.
}
\]

This has exactly the same mathematical structure as the Hammer--Rosen
diffusion equation.

\section{Generalized CTE Form Including Radiation Energy}

If we still assume
\[
T_R = T_m = T,
\]
but do not neglect the radiation energy \(aT^4\), then the conserved energy is
\[
U(T) = A_mT^\beta + aT^4 .
\]

The diffusion variable is
\[
\Phi(T)
=
\int_0^T
D_0T'^\alpha
\frac{d(aT'^4)}{dT'}\,dT' .
\]

Since
\[
\frac{d(aT'^4)}{dT'} = 4aT'^3,
\]
we get
\[
\Phi(T)
=
4aD_0\int_0^T T'^{\alpha+3}\,dT',
\]
and therefore
\[
\boxed{
\Phi(T)
=
\frac{4aD_0}{4+\alpha}T^{4+\alpha}.
}
\]

Let
\[
\Phi_s(t)=\Phi(T_s(t)).
\]
Following the Hammer--Rosen arbitrary-material-property formulation, the heat
front may be written as
\[
\boxed{
x_F^2(t)
\simeq
\frac{2}{L_s(t)}
\int_0^t \Phi_s(t')\,dt',
}
\]
where
\[
\boxed{
L_s(t)
=
\frac{2}{\Phi_s(t)}
\int_0^{\Phi_s(t)}
U(\Phi)
\left(
1-\frac{\Phi}{\Phi_s(t)}
\right)d\Phi .
}
\]

The corresponding total areal energy is
\[
\boxed{
\mathcal{E}(t)
\simeq
x_F(t)
\frac{1}{\Phi_s(t)}
\int_0^{\Phi_s(t)}
U(\Phi)\,d\Phi .
}
\]

For the power-law form
\[
U(T)=A_mT^\beta+aT^4,
\]
define
\[
\epsilon_m = \frac{\beta}{4+\alpha},
\qquad
\epsilon_r = \frac{4}{4+\alpha}.
\]
Then
\[
\boxed{
L_s
=
2\left[
\frac{A_mT_s^\beta}
{(1+\epsilon_m)(2+\epsilon_m)}
+
\frac{aT_s^4}
{(1+\epsilon_r)(2+\epsilon_r)}
\right].
}
\]

Similarly,
\[
\boxed{
\mathcal{E}(t)
\simeq
x_F(t)
\left[
\frac{A_mT_s^\beta}{1+\epsilon_m}
+
\frac{aT_s^4}{1+\epsilon_r}
\right].
}
\]

\section{True Two-Temperature Case}

For the full gray system, where
\[
T_R \neq T_m,
\]
one can still derive an exact energy-conservation equation, but the problem is
not closed by energy conservation alone.

Introduce the normalized variables
\[
y=\frac{x}{x_F(t)},
\qquad
E(x,t)=aT_s^4(t)R(y,t),
\qquad
T_m(x,t)=T_s(t)M(y,t).
\]

Then the total energy is
\[
\mathcal{E}(t)
=
x_F(t)
\left[
aT_s^4(t)I_R(t)
+
A_mT_s^\beta(t)I_M(t)
\right],
\]
where
\[
I_R(t)=\int_0^1 R(y,t)\,dy,
\qquad
I_M(t)=\int_0^1 M^\beta(y,t)\,dy .
\]

The incoming radiation flux is
\[
F_{\rm in}(t)
=
\frac{D_0aT_s^{4+\alpha}(t)}{x_F(t)}
A_R(t),
\]
where
\[
A_R(t)
=
-
\left.
M^\alpha(y,t)
\frac{\partial R}{\partial y}
\right|_{y=0}.
\]

Therefore, the exact moving-front energy equation is
\[
\boxed{
\frac{d}{dt}
\left\{
x_F(t)
\left[
aT_s^4(t)I_R(t)
+
A_mT_s^\beta(t)I_M(t)
\right]
\right\}
=
\frac{D_0aT_s^{4+\alpha}(t)}{x_F(t)}
A_R(t).
}
\]

This equation is exact, but it is not closed because the quantities
\[
I_R(t), \qquad I_M(t), \qquad A_R(t)
\]
depend on the unknown two-temperature profiles \(R(y,t)\) and \(M(y,t)\).

Thus, in the true two-temperature case, a closed analytic expression for
\(x_F(t)\) analogous to the Hammer--Rosen result generally requires solving a
new coupled eigenvalue or self-similar problem. Only in the near-equilibrium
limit,
\[
T_R \simeq T_m,
\]
does the system reduce to a Hammer--Rosen-type single-temperature diffusion
problem.

\section{Henyey-Type Two-Temperature Ansatz}

We now ask whether the two-temperature gray system can be reduced by assuming
Henyey-like spatial profiles for the radiation and material fields. The gray
equations are
\begin{align}
\frac{\partial E}{\partial t}
-
\frac{\partial}{\partial x}
\left(
D_0T_m^\alpha
\frac{\partial E}{\partial x}
\right)
&=
S_0T_m^{-\alpha}
\left(aT_m^4-E\right),
\\
A_m\frac{\partial T_m^\beta}{\partial t}
&=
S_0T_m^{-\alpha}
\left(E-aT_m^4\right),
\end{align}
where
\[
D_0=\frac{cg}{3}\rho^{-1-\lambda},
\qquad
S_0=\frac{c}{g}\rho^{1+\lambda},
\qquad
A_m=f\rho^{1-\mu}.
\]

We introduce the moving-front coordinate
\[
y=\frac{x}{x_F(t)},
\qquad 0\leq y\leq 1.
\]

We allow the radiation and material temperatures to differ:
\[
T_R(x,t)=T_{R,s}(t)\,\mathcal{R}(y),
\qquad
T_m(x,t)=T_s(t)M(y),
\]
and define
\[
E(x,t)=E_s(t)R(y),
\qquad
R(y)\equiv \mathcal{R}^4(y),
\]
where
\[
E_s(t)=aT_{R,s}^4(t).
\]

A Henyey-like profile may be written as
\[
R(y)=(1-y)(1+\chi_R y),
\]
and, for the material energy profile,
\[
M^\beta(y)=(1-y)(1+\chi_M y).
\]

The constants \(\chi_R\) and \(\chi_M\) describe the curvature of the radiation
and material profiles. The Hammer--Rosen first-order Henyey profile has the
same structure:
\[
U(y)\sim (1-y)(1+\chi y).
\]

The boundary conditions are
\[
R(0)=1,
\qquad
M(0)=1,
\qquad
R(1)=0,
\qquad
M(1)=0.
\]

\section{Profile Integrals}

Define the radiation and material energy-shape integrals
\[
I_R=\int_0^1 R(y)\,dy,
\qquad
I_M=\int_0^1 M^\beta(y)\,dy.
\]

For the Henyey-like profiles above,
\[
I_R=\int_0^1 (1-y)(1+\chi_R y)\,dy
=
\frac{1}{2}+\frac{\chi_R}{6},
\]
and similarly
\[
I_M=
\frac{1}{2}+\frac{\chi_M}{6}.
\]

The surface-gradient coefficient is
\[
A_R=
-\left.
M^\alpha(y)\frac{dR}{dy}
\right|_{y=0}.
\]

Since \(M(0)=1\),
\[
A_R=-R'(0).
\]

For
\[
R(y)=(1-y)(1+\chi_Ry),
\]
we have
\[
R'(y)=\chi_R-1-2\chi_R y,
\]
and therefore
\[
R'(0)=\chi_R-1.
\]

Hence
\[
\boxed{
A_R=1-\chi_R.
}
\]

The incoming radiative flux is therefore
\[
F_{\rm in}
=
-
\left.
D_0T_m^\alpha
\frac{\partial E}{\partial x}
\right|_{x=0}
=
\frac{D_0T_s^\alpha E_s}{x_F}A_R.
\]

\section{Total Energy and Heat-Front Equation}

The total areal energy inside the heated region is
\[
\mathcal{E}(t)
=
\int_0^{x_F(t)}
\left[
E(x,t)+A_mT_m^\beta(x,t)
\right]dx.
\]

Using the assumed profiles,
\[
\mathcal{E}(t)
=
x_F(t)
\left[
E_s(t)I_R
+
A_mT_s^\beta(t)I_M
\right].
\]

Define
\[
B(t)=E_s(t)I_R+A_mT_s^\beta(t)I_M.
\]

Then
\[
\mathcal{E}(t)=x_F(t)B(t).
\]

Energy conservation gives
\[
\frac{d}{dt}\left[x_F(t)B(t)\right]
=
\frac{D_0T_s^\alpha(t)E_s(t)}{x_F(t)}A_R.
\]

Define
\[
K(t)=D_0A_R T_s^\alpha(t)E_s(t).
\]

Then the moving-front equation becomes
\[
\boxed{
\frac{d}{dt}\left[x_F(t)B(t)\right]
=
\frac{K(t)}{x_F(t)}.
}
\]

Multiplying by \(x_F\), we get
\[
B x_F\dot{x}_F+x_F^2\dot{B}=K.
\]

Since
\[
x_F\dot{x}_F=\frac{1}{2}\frac{d x_F^2}{dt},
\]
we obtain
\[
\frac{B}{2}\frac{d x_F^2}{dt}
+
\dot{B}x_F^2
=
K.
\]

Let
\[
X(t)=x_F^2(t).
\]

Then
\[
\frac{B}{2}\dot{X}+\dot{B}X=K.
\]

Dividing by \(B/2\),
\[
\dot{X}
+
2\frac{\dot{B}}{B}X
=
\frac{2K}{B}.
\]

The integrating factor is
\[
\exp\left(
\int 2\frac{\dot{B}}{B}\,dt
\right)
=
B^2.
\]

Therefore,
\[
\frac{d}{dt}
\left[
B^2(t)X(t)
\right]
=
2K(t)B(t).
\]

Since \(X=x_F^2\), the heat-front position is
\[
\boxed{
x_F^2(t)
=
\frac{2}{B^2(t)}
\int_0^t K(t')B(t')\,dt'.
}
\]

That is,
\[
\boxed{
x_F^2(t)
=
\frac{2}{\left[
E_s(t)I_R+A_mT_s^\beta(t)I_M
\right]^2}
\int_0^t
D_0A_R T_s^\alpha(t')E_s(t')
\left[
E_s(t')I_R+A_mT_s^\beta(t')I_M
\right]
dt'.
}
\]

This is the two-temperature Henyey-profile analogue of the Hammer--Rosen
front-position formula.

\section{Total Energy}

Since
\[
\mathcal{E}(t)=x_F(t)B(t),
\]
we immediately obtain
\[
\boxed{
\mathcal{E}^2(t)
=
2
\int_0^t
K(t')B(t')\,dt'.
}
\]

Explicitly,
\[
\boxed{
\mathcal{E}^2(t)
=
2
\int_0^t
D_0A_R T_s^\alpha(t')E_s(t')
\left[
E_s(t')I_R+A_mT_s^\beta(t')I_M
\right]
dt'.
}
\]

Therefore,
\[
\boxed{
\mathcal{E}(t)
=
\left\{
2
\int_0^t
D_0A_R T_s^\alpha(t')E_s(t')
\left[
E_s(t')I_R+A_mT_s^\beta(t')I_M
\right]
dt'
\right\}^{1/2}.
}
\]

\section{Closure for the Radiation--Material Temperature Difference}

The previous result gives a front equation, but the two-temperature problem is
not fully closed unless \(E_s(t)\) and \(T_s(t)\) are known. To close the model,
one can integrate the material energy equation over the heated region.

The material energy is
\[
\mathcal{E}_m(t)
=
x_F(t)A_mT_s^\beta(t)I_M.
\]

The volume-integrated matter-radiation exchange is
\[
Q_{\rm abs}(t)
=
\int_0^{x_F}
S_0T_m^{-\alpha}
\left(
E-aT_m^4
\right)dx.
\]

Using the profile ansatz,
\[
Q_{\rm abs}(t)
=
x_F S_0T_s^{-\alpha}
\left[
E_s J_R
-
aT_s^4 J_M
\right],
\]
where
\[
J_R=
\int_0^1 M^{-\alpha}(y)R(y)\,dy,
\]
and
\[
J_M=
\int_0^1 M^{4-\alpha}(y)\,dy.
\]

Therefore the material-energy equation becomes
\[
\boxed{
\frac{d}{dt}
\left[
x_F A_mT_s^\beta I_M
\right]
=
x_F S_0T_s^{-\alpha}
\left[
E_s J_R
-
aT_s^4 J_M
\right].
}
\]

Similarly, the radiation-energy equation is
\[
\boxed{
\frac{d}{dt}
\left[
x_F E_s I_R
\right]
=
\frac{D_0A_RT_s^\alpha E_s}{x_F}
-
x_F S_0T_s^{-\alpha}
\left[
E_s J_R
-
aT_s^4 J_M
\right].
}
\]

Together with
\[
\boxed{
x_F^2(t)
=
\frac{2}{B^2(t)}
\int_0^t K(t')B(t')\,dt',
}
\]
these equations form a closed reduced two-temperature model, provided the
shape parameters \(\chi_R\) and \(\chi_M\) are prescribed or fitted.

\section{Equilibrium Limit}

In the complete thermodynamic equilibrium limit,
\[
T_R=T_m=T_s,
\]
so
\[
E_s=aT_s^4.
\]

If, in addition, the radiation and material profiles have the same shape, then
\[
R(y)\approx M^4(y).
\]

The total energy coefficient becomes
\[
B(t)
=
aT_s^4 I_R
+
A_mT_s^\beta I_M.
\]

If material energy dominates radiation energy,
\[
A_mT_s^\beta I_M \gg aT_s^4 I_R,
\]
then
\[
B(t)\approx A_mT_s^\beta I_M.
\]

Also,
\[
K(t)=D_0A_RaT_s^{4+\alpha}.
\]

Therefore,
\[
x_F^2(t)
\approx
\frac{2D_0A_Ra}{A_m I_M}
T_s^{-2\beta}(t)
\int_0^t
T_s^{4+\alpha+\beta}(t')\,dt'.
\]

This has the same diffusion-wave character as the Hammer--Rosen result, but
with constants determined by the chosen Henyey-like profile.
\end{document}